\chapter{Zeitdiskrete Systeme im Zeitbereich}\label{sec:zeitdiskrete-systeme-im-zeitbereich}

Im zeitkontinuierlichen Bereich werden dynamische Systeme mit Differentialgleichungen beschrieben. Für zeitdiskrete Systeme geht die Differentiation in einen Differenzenquotienten über. Daraus resultiert die Beschreibung zeitdiskreter Systeme mit Differenzengleichungen.

Zu Beginn des Kapitels werden einige typische zeitdiskrete Systeme beschrieben und die Systemantworten mit numerischen Algorithmen bestimmt. Als Beispiele werden ein rekursives Filter, ein System mit einer gleitenden Mittelwertrechnung und eine Switched-Capacitor-Schaltung diskutiert. Für die Systeme werden die jeweiligen Differenzengleichungen aufgestellt, diskutiert und numerisch gelöst.

Lineare, zeitinvariante Systeme haben im zeitkontinuierlichen Bereich eine besonders anschauliche Lösung und Interpretation ermöglicht. Der Begriff des linearen, zeitinvarianten Systems existiert auch für zeitdiskrete Systeme. Im zweiten Abschnitt dieses Kapitels werden grundlegende Systemeigenschaften erläutert und Verfahren zum Nachweis der unterschiedlichen Eigenschaften vorgestellt.

Numerische Lösungen von Differenzengleichungen eignen sich nicht für eine analytische Bestimmung von Systemeigenschaften. Deshalb ist es erforderlich, analytische Lösungsmethoden zu diskutieren. Im dritten Abschnitt wird die analytische Berechnung der Systemantwort im Zeitbereich behandelt. Sie kann über eine sogenannte 4-Schritt-Methode bestimmt werden.

Für lineare, zeitinvariante Systeme gilt das Superpositionsprinzip. Danach ergibt sich die Systemantwort bei Anregung mit einer Überlagerung von Eingangssignalen aus derselben Überlagerung der zugehörigen Ausgangssignale. Das Superpositionsprinzip führt zur zeitdiskreten Faltung. Superposition und Faltung entsprechen weitgehend den Methoden für zeitkontinuierliche Systeme.

Verfahren und Eigenschaften zeitdiskreter Systeme werden an Beispielen erläutert. An einem Projekt wird aufgezeigt, wie zeitdiskrete Signalverarbeitung mit Switched-Capacitor-Schaltungen realisiert werden kann und welche Vorteile sich aus der Realisierung ergeben.

\section{Beschreibung zeitdiskreter Systeme mit Differenzengleichungen}\label{sec:4.1}

Zeitdiskrete Systeme ergeben sich aus unterschiedlichen Aufgabenstellungen. Zum einen existieren Aufgabenstellungen, die von ihrer Natur aus zeitdiskret sind. Beispiele dafür kommen aus unterschiedlichen Anwendungsbereichen, zum Beispiel werden Algorithmen zur digitalen Signalverarbeitung zeitdiskret beschrieben. Aber auch Schaltungen, die in Switched-Capacitor-Technik aufgebaut werden, können zeitdiskret beschrieben werden.

Zum anderen wird vielfach versucht, analoge Systeme digital nachzubilden, um ihr Verhalten zu emulieren. Lineare Differentialgleichungen N-ter Ordnung mit konstanten Koeffizienten beschreiben zeitkontinuierliche LTI-Systeme. Bei abgetasteten oder zeitdiskreten Systemen kann keine Ableitung bestimmt werden. Aus diesem Grund muss eine andere Beschreibungsform für Systeme gefunden werden. In diesem Abschnitt werden unterschiedliche Methoden vorgestellt, LTI-Systeme zeitdiskret zu beschreiben.

\subsection{Beispiele für zeitdiskrete Systeme}\label{sec:beispiele-fuxfcr-zeitdiskrete-systeme}

Algorithmen zur digitalen Signalverarbeitung arbeiten typischerweise in einem festen Takt mit einer festen Abtastzeit. Der aktuelle Wert des Ausgangssignals wird als Funktion des aktuellen Eingangssignals und vergangenen Werten von Ein- und Ausgangssignal gebildet. Das Grundprinzip wird an Beispielen illustriert.

\begin{example}{Rückgekoppeltes System}
\label{beispiel-ruxfcckgekoppeltes-system}

Bei vielen Geräten wird versucht, Filter mit wenig Rechenzeit und wenig Speicherelementen zu implementieren. Eines der einfachsten Filter ergibt sich aus der Gleichung
\begin{align}
y[ k] &= (1 - GF)  u[ k] + GF  y[ k - 1] \label{eq:4.1}
\end{align}
Dabei gilt für den Gedächtnisfaktor GF die Bedingung 0 \textless{} GF \textless{} 1. Das aktuelle Ausgangssignal ergibt sich aus dem aktuellen Eingangssignal und dem um einen Takt zurückliegenden Ausgangssignal. Für kleine Werte von GF folgt das Ausgangssignal dem Eingangssignal, für große Werte ändert sich das Ausgangsignal nur sehr langsam. Bild 4.1 zeigt diesen Zusammenhang für einen Sprung am Eingang.

\begin{fixedfigure}
  \includegraphics[alt={Ausgangssignal eines zeitdiskreten Tiefpass-Filters für einen Sprung am Eingang}]{new media/4_1_ausgangssignal_eines_zeitdiskret.png}
  \captionof{figure}{Ausgangssignal eines zeitdiskreten Tiefpass-Filters für einen Sprung am Eingang}
\end{fixedfigure}

Die Systemantworten erinnern an das Ausgangssignal eines Tiefpasses. Es wird sich zeigen, dass diese rekursive Form die digitale Approximation eines Tiefpasses erster Ordnung ist.

\end{example}

\begin{example}{Gleitender Mittelwert}
\label{beispiel-gleitender-mittelwert}

Zur Verringerung von Störeffekten kann der Mittelwert einiger aufeinanderfolgender Folgenwerte gebildet werden. Dieses Vorgehen wird als gleitende Mittelwertbildung bezeichnet. Zum Beispiel wird über die Gleichung
\begin{align}
y[ k] &= \frac{1}{5}  \left( u[ k] + u[ k - 1] + u[ k - 2] + u[ k - 3] + u[ k - 4] \right) \label{eq:4.2}
\end{align}
ein gleitendes Mittelwertfilter mit 5 Abtastwerten beschrieben. Die Ausgangsfolge ist in diesem Fall nur von aktuellen und vergangenen Werten des Eingangssignals abhängig. Durch die Mittelwertbildung weist die Ausgangsfolge \(y[k]\) einerseits eine geringere Variation auf als die Eingangsfolge \(u[k]\). Durch die Bildung des Mittelwertes von fünf Abtastwerten erreicht das Signal andererseits später seinen Endwert. Beide Effekte sind in Bild \ref{fig:4.2} verdeutlicht.

\begin{fixedfigure}
  \includegraphics[alt={Ausgangssignal eines zeitdiskreten Systems mit gleitender Mittelung}]{new media/4_2_ausgangssignal_eines_zeitdiskret.png}
  \captionof{figure}{Ausgangssignal eines zeitdiskreten Systems mit gleitender Mittelung}
  \label{fig:4.2}
\end{fixedfigure}

\end{example}
\begin{application}{Zeitdiskrete Systeme}
  Im Online-Portal \emph{Systemtheorie Online} verdeutlicht die Applikation \emph{Zeitdiskrete Systeme} die Wirkung eines rekursiven Tiefpass-Filters und eines Mittelwert-Filters auf ein verrauschtes Signal.
  \tcblower
  \url{h-ka.de/eit/digitale-signalverarbeitung/übungen/kapitel2}
\end{application}
Die beiden vorangegangenen Beispiele sind Algorithmen, die zum Beispiel als Programm in einem Controller implementiert werden können. Ein weiteres Anwendungsfeld für zeitdiskrete Systeme sind Switched-Capacitor-Schaltungen.

\begin{example}{Switched-Capacitor-Schaltung}
\label{beispiel-switched-capacitor-schaltung}

Bei der Integration von CMOS-Schaltungen ist es aufwendig und teuer, präzise und temperaturstabile Widerstände herzustellen. Die Herstellung von stabilen Kondensatoren ist dagegen vergleichsweise günstig. Deshalb werden in der CMOS-Technik Switched-Capacitor-Schaltungen verwendet. Als Beispiel zeigt Bild \ref{fig:4.3} einen invertierenden Integrierer in Switched-Capacitor-Technik.

\begin{fixedfigure}
  \includegraphics[alt={Invertierender Integrierer in Switched-Capacitor-Technik}]{new media/4_3_invertierender_integrierer_in_sw.png}
  \captionof{figure}{Invertierender Integrierer in Switched-Capacitor-Technik}
  \label{fig:4.3}
\end{fixedfigure}

Zu Beginn jeden Taktes wird der Kondensator \(C_{E}\) aufgeladen. Nach Umschalten des Schalters wird der Kondensator \(C_{E}\) mit dem Knoten A des Operationsverstärkers verbunden, der auf Ground-Potential liegt. Der Kondensator wird dabei entladen. Da der Operationsverstärker einen sehr großen Innenwiderstand besitzt, fließt die Ladung nicht über den Operationsverstärker ab, sondern zum Kondensator \(C_{A}\). Bei jedem Takt wird also dem Integrierer eine Ladung der Größe \(q_{E}\){[}k{]} zugeführt.
\begin{align}
q_{E}[ k] &= u_{E}[ k]  C_{E} \label{eq:4.3}
\end{align}
Die aktuelle Ladung im Kondensator \(C_{A}\) beträgt damit
\begin{align}
q_{A}[ k] &= u_{E}[ k]  C_{E} + q_{A}[ k - 1] = - u_{A}[ k]  C_{A} \label{eq:4.4}
\end{align}
Die aktuelle Ladung \(q_{A}\){[}k{]} ergibt sich aus der Ladung im Takt zuvor und der dazukommenden Ladung. Die Ladung wird damit aufsummiert. Das aktuelle Ausgangssignal ergibt sich zu
\begin{align}
u_{A}[ k] &= - \frac{1}{C_{A}}  q_{A}[ k] = - \frac{1}{C_{A}}  \left( u_{E}[ k]  C_{E} + q_{A}[ k - 1] \right) \label{eq:4.5}
\end{align}
Wegen des Minuszeichens ist das System ein invertierender Summierer.

\end{example}

Alle hier behandelten Beispiele führen auf eine Differenzengleichung, die in allgemeiner Form geschrieben werden kann als
\begin{align}
c_{0}  y[ k] + c_{1}  y[ k - 1] + ... + c_{n}  y[ k - N] &= d_{0}  u[ k] + d_{1}  u[ k - 1] + ... + d_{L}  u[ k - L] \label{eq:4.6}
\end{align}
Die Lösung dieser Differenzengleichung ist Gegenstand des Abschnitts \hyperref[_Ref369028047]{0}.

\subsection{Zeitdiskrete Approximation zeitkontinuierlicher Systeme}\label{sec:zeitdiskrete-approximation-zeitkontinuierlicher-systeme}

Das Verhalten der Ein- und Ausgangsgrößen eines realisierbaren, zeitkontinuierlichen LTI-Systems wird durch eine lineare Differentialgleichung N-ter Ordnung mit konstanten Koeffizienten beschrieben. Nach den Ausführungen in Teil A dieser Buchreihe kann das System über Integrierer dargestellt werden. Beim Übergang von zeitkontinuierlichen zu zeitdiskreten Systemen kann die Integration nur approximiert werden. Eine Realisierungsmöglichkeit ist die Approximation mit einem Trapez. Bild \ref{fig:4.4} verdeutlicht die Integration und Approximation der Integration grafisch.

\begin{figure}
  \centering
  \includegraphics[alt={Approximation des Integrals eines zeitkontinuierlichen Signals durch ein Trapez}]{new media/4_4_approximation_des_integrals_eine.png}
  \caption{Approximation des Integrals eines zeitkontinuierlichen Signals durch ein Trapez}
  \label{fig:4.4}
\end{figure}

Aus einem Vergleich der Flächen unter der Funktion \(u(t)\) ergibt sich der Ausdruck
\begin{align}
\int_{(k - 1)  T_{A}}^{k  T_{A}}{u(t)}\ dt &= \frac{1}{2}  \left( u[ k] + u[ k - 1] \right)  T_{A} \label{eq:4.7}
\end{align}
Die Approximation der Integration führt zu der gewichteten Summe aufeinanderfolgender Abtastwerte, die mit dem Faktor \(T_{A}\) multipliziert wird. Für Integrale höherer Ordnung kann entsprechend verfahren werden. Es ergibt sich eine Summe von Folgenwerten, die jeweils um einen Abtastwert verschoben sind. Allgemein ist das aktuelle Ausgangssignal eine Funktion der vergangenen Ein- und Ausgangswerte und des aktuellen Eingangssignals.
\begin{align}
y[ k] &= f\left( y[ k - 1],y[ k - 2],...,u[ k],u[ k - 1],u[ k - 2],... \right) \label{eq:4.8}
\end{align}
Die Beschreibung linearer, zeitinvarianter Systeme führt zu linearen Differenzengleichungen mit konstanten Koeffizienten, die allgemein dargestellt werden können als
\begin{align}
\sum\limits_{n = 0}^{N}{c_{n}  y[ k - n]} &= \sum\limits_{l = 0}^{L}{d_{l}  u[ k - l]} \label{eq:4.9}
\end{align}
beziehungsweise mit \(c_{0}\) = 1
\begin{align}
y[ k] &= \sum\limits_{l = 0}^{L}{d_{l}  u[ k - l]} - \sum\limits_{n = 1}^{N}{c_{n}  y[ k - n]} \label{eq:4.10}
\end{align}
Die Güte der Approximation verbessert sich mit sinkender Abtastzeit \(T_{A}\). Allerdings ist die Abtastzeit aus der Anwendung vorgegeben. Auch der Signalverlauf hat einen Einfluss auf die Approximationsgüte. Nichtlineare Kurvenverläufe und hochfrequente Signalanteile wirken sich negativ auf die Approximationsgüte aus.

\begin{example}{RC-Tiefpass}
\label{beispiel-rc-tiefpass}

Bild \ref{fig:3.1} zeigt einen RC-Tiefpass mit Spannungsquelle \(u_{E}\), Widerstand R und Kapazität C.

\begin{fixedfigure}
  \includegraphics[alt={RC-Tiefpass}]{new media/4_5_rctiefpass.pdf}
  \captionof{figure}{RC-Tiefpass}
\end{fixedfigure}

Die zugehörige lineare Differentialgleichung für die Ausgangsspannung \(u_{A}(t)\) lautet:
\begin{align}
u_{A}(t) + R  C  \frac{du_{A}}{dt} &= u_{E}(t) \label{eq:4.11}
\end{align}
Eine Integration der Differentialgleichung führt bei zeitdiskreten Integrationsgrenzen zu
\begin{align}
\int_{(k - 1)  T_{A}}^{k  T_{A}}{u_{A}(t)\ }dt + R  C  \left( u_{A}\left( k  T_{A} \right) - u_{A}\left( (k - 1)  T_{A} \right) \right) &= \int_{(k - 1)  T_{A}}^{k  T_{A}}{u_{A}(t)\ }dt \label{eq:4.12}
\end{align}
Mit den Signalen
\begin{align}
u_{E}\left( k  T_{A} \right) &= u_{E}[ k] \label{eq:4.13}
\end{align}
und
\begin{align}
u_{A}\left( k  T_{A} \right) &= u_{A}[ k] \label{eq:4.14}
\end{align}
und einer Approximation der Integrale über die Trapezregel ergibt sich
\begin{align}
\frac{1}{2}  \left( u_{A}[ k] + u_{A}[ k - 1] \right)  T_{A} + R  C  \left( u_{A}[ k] - u_{A}[ k - 1] \right) &= \frac{1}{2}  \left( u_{E}[ k] + u_{E}[ k - 1] \right)  T_{A} \label{eq:4.15}
\end{align}
Zur Berechnung des aktuellen Ausgangssignals wird die Gleichung ausmultipliziert und nach \(u_{A}\){[}k{]} aufgelöst.
\begin{align}
u_{A}[ k] &= \frac{T_{A}  u_{E}[ k] + T_{A}  u_{E}[ k - 1] - T_{A}  u_{A}[ k - 1] + 2  R  C  u_{A}[ k - 1]}{T_{A} + 2  R  C} \label{eq:4.16}
\end{align}
Die aktuelle Ausgangsspannung \(u_{A}\){[}k{]} berechnet sich aus dem aktuellen und einem vergangenen Wert der Eingangsspannung \(u_{E}\){[}k{]} und einem vergangenen Wert der Ausgangsspannung \(u_{A}\){[}k{]}. Bild \ref{fig:4.6} zeigt den Vergleich zwischen dem zeitkontinuierlichen und dem zeitdiskreten System bei unterschiedlichen Abtastzeiten.

\begin{fixedfigure}
  \includegraphics[alt={Vergleich zwischen zeitkontinuierlicher und zeitdiskreter Realisierung eines Tiefpasses (RC = \SI{100}{\micro\second})}]{new media/4_6_vergleich_zwischen_zeitkontinuie.png}
  \captionof{figure}{Vergleich zwischen zeitkontinuierlicher und zeitdiskreter Realisierung eines Tiefpasses (RC = \SI{100}{\micro\second})}
  \label{fig:4.6}
\end{fixedfigure}

Die Abtastzeit von \(T_{A} = \SI{10}{\micro\second}\) ist deutlich kleiner als die Zeitkonstante des Systems mit \(T = \SI{100}{\micro\second}\). Aus diesem Grund ist die Approximationsgüte im links dargestellten Fall hoch. Die im rechten Bildteil diskutierte Abtastzeit von \(T_{A} = \SI{100}{\micro\second}\) entspricht der Zeitkonstante des Systems. Aus diesem Grund ist die Approximation der Ableitung mit einem Differenzenquotienten unzureichend, die Approximationsgüte ist deutlich schlechter als bei einer Abtastzeit von \(T_{A} = \SI{10}{\micro\second}\). Das Beispiel zeigt, dass zeitkontinuierliche Systeme zeitdiskret nachgebildet werden können. Die lineare Differentialgleichung erster Ordnung
\begin{align}
u_{A}(t) + R  C  \frac{du_{A}}{dt} &= u_{E}(t) \label{eq:4.17}
\end{align}
geht dabei über in eine lineare Differenzengleichung erster Ordnung.
\begin{align}
\left( T_{A} + 2  R  C \right)  u_{A}[ k] + \left( T_{A} - 2  R  C \right)  u_{A}[ k - 1] &= T_{A}  u_{E}[ k] + T_{A}  u_{E}[ k - 1] \label{eq:4.18}
\end{align}
\end{example}

Die hier diskutierten Beispiele führen auf lineare Differenzengleichungen mit konstanten Koeffizienten.
\begin{align}
c_{0}  y[ k] + c_{1}  y[ k - 1] + ... + c_{n}  y[ k - N] &= d_{0}  u[ k] + d_{1}  u[ k - 1] + ... + d_{L}  u[ k - L] \label{eq:4.19}
\end{align}
Im Folgenden werden die Eigenschaften von Systemen diskutiert, die sich mit dieser Art von Differenzengleichungen beschreiben lassen.

\section{Grundlegende Systemeigenschaften}\label{sec:4.2}

Im Abschnitt \ref{sec:4.1} werden unterschiedliche Systeme beschrieben. Die Beschreibung führt zu linearen Differenzengleichungen mit konstanten Koeffizienten. In diesem Abschnitt werden grundlegende Eigenschaften von Systemen diskutiert. Sie entsprechen den Eigenschaften linearer, zeitinvarianter Systeme im zeitkontinuierlichen Bereich. Es wird sich zeigen, dass lineare Differenzengleichungen mit konstanten Koeffizienten zeitdiskrete, lineare, zeitinvariante Systeme (LTI-Systeme) beschreiben.

\subsection{Linearität}\label{sec:linearituxe4t}

Für den Linearitätsnachweis eines Systems müssen die Systemantworten \(y_{1}\){[}k{]} und \(y_{2}\){[}k{]} auf die Eingangsfolgen \(u_{1}\){[}k{]} und \(u_{2}\){[}k{]} bekannt sein. Ein System ist linear, wenn es auf eine Linearkombination von Eingangssignalen
\begin{align}
u[ k] &= \nu_{1}  u_{1}[ k] + \nu_{2}  u_{2}[ k] \label{eq:4.20}
\end{align}
mit derselben Linearkombination der entsprechenden Ausgangssignale reagiert.
\begin{align}
y[ k] &= \nu_{1}  y_{1}[ k] + \nu_{2}  y_{2}[ k] \label{eq:4.21}
\end{align}
Der Nachweis erfolgt über Einsetzen der Gleichungen in die Differenzengleichung.

\begin{example}{Linearität eines rückgekoppelten Systems}
\label{beispiel-linearituxe4t-eines-ruxfcckgekoppelten-systems}

Das Filter mit der Differenzengleichung
\begin{align}
y[ k] &= (1 - GF)  u[ k] + GF  y[ k - 1] \label{eq:4.22}
\end{align}
soll auf Linearität untersucht werden. Die Systemantworten \(y_{1}\){[}k{]} und \(y_{2}\){[}k{]} berechnen sich mit der Differenzengleichung zu
\begin{align}
y_{1}[ k] &= (1 - GF)  u_{1}[ k] + GF  y_{1}[ k - 1] \label{eq:4.23}
\end{align}
beziehungsweise
\begin{align}
y_{2}[ k] &= (1 - GF)  u_{2}[ k] + GF  y_{2}[ k - 1] \label{eq:4.24}
\end{align}
Wird das System mit der oben beschriebenen Linearkombination angeregt, ergibt sich das Ausgangssignal \(y[k]\) aus derselben Linearkombination wie die Eingangssignale.
\begin{align}
{y\left[k\right]} &=  {\left(1-GF\right) u\left[k\right]+GF y\left[k-1\right]} \\
{} &=  {\left(1-GF\right) \left(\nu _{1}  u_{1} \left[k\right]+\nu _{2}  u_{2} \left[k\right]\right)+GF \left(\nu _{1}  y_{1} \left[k-1\right]+\nu _{2}  y_{2} \left[k-1\right]\right)} \\
{} &=  {\left(1-GF\right) \nu _{1}  u_{1} \left[k\right]+\left(1-GF\right) \nu _{2}  u_{2} \left[k\right]+GF \nu _{1}  y_{1} \left[k-1\right]+GF \nu _{2}  y_{2} \left[k-1\right]} \\
{} &=  {\nu _{1}  \left(1-GF\right) u_{1} \left[k\right]+\nu _{1}  GF y_{1} \left[k-1\right]+\nu _{2}  \left(1-GF\right) u_{2} \left[k\right]+\nu _{2}  GF y_{2} \left[k-1\right]} \\
{} &=  {\nu _{1}  y_{1} \left[k\right]+\nu _{2}  y_{2} \left[k\right]}
 \label{eq:4.25}
\end{align}
Bild \ref{fig:4.7} zeigt Ein- und Ausgangssignale eines linearen Systems, das mit den Signalen \(u_{1}\){[}k{]}, \(u_{2}\){[}k{]} und \(u[k]\) = \(u_{1}\){[}k{]} + \(u_{2}\){[}k{]} angeregt wird.

\begin{fixedfigure}
  \includegraphics[alt={Reaktion eines linearen Systems auf die Anregung mit einer Linearkombination von Signalen}]{new media/4_7_reaktion_eines_linearen_systems.png}
  \captionof{figure}{Reaktion eines linearen Systems auf die Anregung mit einer Linearkombination von Signalen}
  \label{fig:4.7}
\end{fixedfigure}

Auch das Ausgangssignal \(y[k]\) setzt sich aus der Summe der Ausgangsignale \(y_{1}\){[}k{]} und \(y_{2}\){[}k{]} zusammen.

\end{example}

Die Linearität des Systems kann auch daran abgelesen werden, dass alle Folgen nur in linearen Summen auftreten. Keine der Folgen besitzt eine Potenz ungleich 1 oder tritt als Produkt von Folgen auf.

\subsection{Zeitinvarianz}\label{sec:zeitinvarianz}

Ein System reagiert auf ein Eingangssignal \(u[k]\) mit einer Systemantwort \(y[k]\). Ist das System zeitinvariant, so reagiert das System auf das verzögerte Eingangssignal u{[}k - \(k_{0}\){]} mit dem Ausgangsignal\\
y{[}k - \(k_{0}\){]}. Zeitinvariante Systeme reagieren also unabhängig vom Startzeitpunkt der Beobachtung auf gleiche Eingangssignale mit gleichen Ausgangssignalen.

\begin{example}{Zeitinvarianz eines rückgekoppelten Systems}
\label{beispiel-zeitinvarianz-eines-ruxfcckgekoppelten-systems}

Das Filter mit der Differenzengleichung
\begin{align}
y[ k] &= (1 - GF)  u[ k] + GF  y[ k - 1] \label{eq:4.26}
\end{align}
soll auf Zeitinvarianz untersucht werden. Dazu werden alle Ausdrücke k durch k -- \(k_{0}\) ersetzt. Es ergibt sich die Differenzengleichung
\begin{align}
y\left[ k - k_{0} \right] &= (1 - GF)  u\left[ k - k_{0} \right] + GF  y\left[ k - k_{0} - 1 \right] \label{eq:4.27}
\end{align}
Wird das Eingangssignal um \(k_{0}\) verschoben, wird auch das Ausgangsignal um \(k_{0}\) verschoben.

\begin{fixedfigure}
  \includegraphics[alt={Reaktion eines zeitinvarianten Systems auf die Anregung mit zeitverschobenen Signalen}]{new media/4_8_reaktion_eines_zeitinvarianten_s.png}
  \captionof{figure}{Reaktion eines zeitinvarianten Systems auf die Anregung mit zeitverschobenen Signalen}
\end{fixedfigure}

\end{example}

Es kann gezeigt werden, dass zeitinvariante Systeme Differenzengleichungen mit konstanten Koeffizienten aufweisen. Ändern sich die Koeffizienten der Differenzengleichung mit dem Folgenindex k, ist das System zeitvariant.

\subsection{Lineare, zeitinvariante Systeme (LTI-Systeme)}\label{sec:4.2.3}

Die Darstellungen in diesem Skript beschränken sich bis auf wenige Ausnahmen auf Systeme, die linear und zeitinvariant sind und als LTI-Systeme bezeichnet werden. Für LTI-Systeme sind besonders anschauliche und einfach zu interpretierende Lösungs- und Interpretationsmethoden im Zeit-, Bild- und Frequenzbereich vorhanden. Sie entsprechen sinngemäß den Methoden zeitkontinuierlicher Systeme.

\begin{figure}
  \centering
  \includegraphics[alt={Bild 4.9: Lineares zeitinvariantes System}]{new media/4_9_bild_49_lineares_zeitinvariantes.pdf}
  \caption{Bild 4.9: Lineares zeitinvariantes System}
\end{figure}

Systeme, die mit einer linearen Differenzengleichung mit konstanten Koeffizienten beschrieben werden können, erfüllen die Bedingungen nach Linearität und Zeitinvarianz. Ausgehend von der Differenzengleichung
\begin{align}
\sum\limits_{n = 0}^{N}{c_{n}  y[ k - n]} &= \sum\limits_{l = 0}^{L}{d_{l}  u[ k - l]} \label{eq:4.28}
\end{align}
werden die Eigenschaften der Linearität und Zeitinvarianz hergeleitet.

\subsubsection{Linearität}\label{sec:linearituxe4t-1}

Ausgangspunkt für den Beweis der Linearität sind zwei Signalkombinationen \(u_{1}\){[}k{]} und \(y_{1}\){[}k{]} sowie \(u_{2}\){[}k{]} und \(y_{2}\){[}k{]}, für die gilt
\begin{align}
\sum\limits_{n = 0}^{N}{c_{n}  y_{1}[ k - n]} &= \sum\limits_{l = 0}^{L}{d_{l}  u_{1}[ k - l]} \label{eq:4.29} \\
\sum\limits_{n = 0}^{N}{c_{n}  y_{2}[ k - n]} &= \sum\limits_{l = 0}^{L}{d_{l}  u_{2}[ k - l]} \label{eq:4.30}
\end{align}
Für eine beliebige Linearkombination
\begin{align}
u[ k] &= \nu_{1}  u_{1}[ k] + \nu_{2}  u_{2}[ k] \label{eq:4.31}
\end{align}
gilt dann die Gleichung
\begin{align}
{\sum\limits _{l=0}^{L}d_{l}  u\left[k-l\right] } &=  {\sum\limits _{l=0}^{L}d_{l}  \nu _{1}  u_{1} \left[k-l\right]+d_{l}  \nu _{2}  u_{2} \left[k-l\right] } \\
{} &=  {\nu _{1}  \sum\limits _{l=0}^{L}d_{l}  u_{1} \left[k-l\right] +\nu _{2}  \sum\limits _{l=0}^{L}d_{l}  u_{2} \left[k-l\right] }
 \label{eq:4.32}
\end{align}
Für jede der beiden Summen kann der Ausdruck aus \eqref{eq:4.29} beziehungsweise \eqref{eq:4.30} eingesetzt werden, und es ergibt sich
\begin{align}
{\sum\limits _{l=0}^{L}d_{l}  u\left[k-l\right] } &=  {\nu _{1}  \sum\limits _{n=0}^{N}c_{n}  y_{1} \left[k-n\right] +\nu _{2}  \sum\limits _{n=0}^{N}c_{n}  y_{2} \left[k-n\right] } \\
{} &=  {\sum\limits _{n=0}^{N}c_{n}  \left(\nu _{1}  y_{1} \left[k-n\right]+\nu _{2}  y_{2} \left[k-n\right]\right) =\sum\limits _{n=0}^{N}c_{n}  y\left[k-n\right] }
 \label{eq:4.33}
\end{align}
Eine Linearkombination von Eingangssignalen führt damit zu der identischen Linearkombination von Ausgangssignalen, sodass das System ein lineares System ist.

\subsubsection{Zeitinvarianz}\label{sec:zeitinvarianz-1}

Wie oben bereits beschrieben sind lineare Differenzengleichungen mit konstanten Koeffizienten stets zeitinvariant. Wird in der Differenzengleichung
\begin{align}
\sum\limits_{n = 0}^{N}{c_{n}  y[ k - n]} &= \sum\limits_{l = 0}^{L}{d_{l}  u[ k - l]} \label{eq:4.34}
\end{align}
der Folgenindex sowohl bei Eingangs- als auch Ausgangssignalen um \(k_{0}\) verschoben, bleibt die Differenzengleichung weiter gültig.
\begin{align}
\sum\limits_{n = 0}^{N}{c_{n}  y\left[ k - k_{0} - n \right]} &= \sum\limits_{l = 0}^{L}{d_{l}  u\left[ k - k_{0} - l \right]} \label{eq:4.35}
\end{align}
\subsection{Kausalität}\label{sec:kausalituxe4t}

Wesentliche Voraussetzung für die Realisierbarkeit eines Systems ist die Forderung, dass die Werte der Ausgangssignale zu einem gegebenen Zeitpunkt nur von Werten der Eingangssignale zu diesem oder einem früheren Zeitpunkt abhängen, nicht jedoch etwa von zukünftigen Werten. Dieses Verhalten wird für zeitkontinuierliche Systeme als die Eigenschaft der Kausalität eines Systems definiert, die auch im diskreten Fall wesentlich ist. Liegt eine Systembeschreibung über eine Differenzengleichung vor, kann die Kausalität direkt bewertet werden.
\begin{align}
\sum\limits_{n = 0}^{N}{c_{n}  y[ k - n]} &= \sum\limits_{l = 0}^{L}{d_{l}  u[ k - l]} \label{eq:4.36}
\end{align}
Ist der Koeffizient \(c_{0}\) ≠ 0, kann ohne Beschränkung der Allgemeinheit die Annahme \(c_{0} = 1\) getroffen werden. Ist das nicht der Fall, wird die Gleichung durch \(c_{0}\) dividiert. Mit dieser Annahme kann die Differenzengleichung nach \(y[k]\) aufgelöst werden, und es ergibt sich
\begin{align}
y[ k] &= \sum\limits_{l = 0}^{L}{d_{l}  u[ k - l]} - \sum\limits_{n = 1}^{N}{c_{n}  y[ k - n]} \label{eq:4.37}
\end{align}
Da alle Indizes l und n größer gleich null sind, ist ein System, das durch eine lineare Differenzengleichung der Form aus Gleichung \eqref{eq:4.37} beschrieben werden kann, ein kausales System. Ist in der Differenzengleichung \eqref{eq:4.36} der Koeffizient \(c_{0} = 0\), kann nicht nach \(y[k]\) aufgelöst werden. Wird nach\\
\(y[k - 1]\) aufgelöst, ergibt sich die Gleichung
\begin{align}
y[ k - 1] &= \frac{1}{c_{1}}  \left( \sum\limits_{l = 0}^{L}{d_{l}  u[ k - l]} - \sum\limits_{n = 2}^{N}{c_{n}  y[ k - n]} \right) \label{eq:4.38}
\end{align}
Damit ist der Ausgangswert \(y[k - 1]\) von dem zukünftigen Eingangswert \(u[k]\) abhängig. Das System ist demnach für \(c_{0} = 0\) nicht kausal.

Durch die Kausalitätsforderung sind einige Systembetrachtungen der digitalen Signalverarbeitung aufwendig. Daher werden im Folgenden aus Vereinfachungsgründen auch nicht-kausale Systeme betrachtet. Diese Systeme können vielfach durch eine Verschiebung kausal und damit realisierbar gemacht werden.

\begin{example}{Kausalität des gleitenden Mittelwertes}
\label{beispiel-kausalituxe4t-des-gleitenden-mittelwertes}

Bei der Beschreibung zeitdiskreter Systeme wird der gleitende Mittelwert vorgestellt. Er hat die Differenzengleichung
\begin{align}
y_{1}[ k] &= \frac{1}{5}  \left( u[ k] + u[ k - 1] + u[ k - 2] + u[ k - 3] + u[ k - 4] \right) \label{eq:4.39}
\end{align}
Weil das Ausgangssignal \(y[k]\) nur von aktuellen und vergangenen Eingangswerten abhängt, ist das System kausal. Es weist aber eine zeitliche Verzögerung auf, die bei der Berechnung von Frequenzgängen noch dargestellt wird. Ein System mit der Differenzengleichung
\begin{align}
y_{2}[ k] &= \frac{1}{5}  \left( u[ k + 2] + u[ k + 1] + u[ k] + u[ k - 1] + u[ k - 2] \right) \label{eq:4.40}
\end{align}
ist nicht mehr kausal, weist aber - wie später gezeigt wird - keine zeitliche Verzögerung auf. Oft werden in der Datenverarbeitung erst alle Messwerte aufgenommen und anschließend ausgewertet. In diesem Fall liegen auch zukünftige Werte vor. Die Kausalität von Systemen ist in diesem Fall nicht zwingend erforderlich. Bild \ref{fig:4.10} stellt die Sprungantwort eines kausalen und eines nicht kausalen Systems zur Berechnung eines gleitenden Mittelwertes dar.

\begin{fixedfigure}
  \includegraphics[alt={Sprungantwort eines kausalen und eines nicht kausalen Systems\\
zur Berechnung eines gleitenden Mittelwertes}]{new media/4_10_sprungantwort_eines_kausalen_un.png}
  \captionof{figure}{Sprungantwort eines kausalen und eines nicht kausalen Systems\\
zur Berechnung eines gleitenden Mittelwertes}
  \label{fig:4.10}
\end{fixedfigure}

Es ist deutlich zu erkennen, dass das nicht kausale System reagiert, bevor der Sprung des Eingangssignals stattgefunden hat, während das kausale Signal erst nach der eigentlichen Anregung reagiert.

\end{example}

\subsection{Stabilität}\label{sec:stabilituxe4t}

Für zeitkontinuierliche Systeme wird in Teil A der Buchreihe folgende Stabilitätsdefinition eingeführt: Ein System ist asymptotisch stabil, wenn es nach einer zeitlich begrenzten Anregung mit endlicher Energie wieder seine Ruheposition erreicht. Es ist grenzstabil, wenn es nach einer zeitlich begrenzten Anregung mit endlicher Energie zu einem konstanten Ausgangswert konvergiert oder mit konstant schwingendem Ausgangssignal reagiert. Ein System ist instabil, wenn es auf eine zeitlich begrenzte Anregung endlicher Energie mit divergierendem Ausgangssignal reagiert. Diese Stabilitätsdefinition wird für zeitdiskrete Systeme übernommen.

Wie bei zeitkontinuierlichen Systemen erfolgt der Nachweis im Zeitbereich über die charakteristische Gleichung (Kapitel \ref{sec:4.3.3}) oder die Impulsantwort (Kapitel \ref{sec:4.4.4}). Alternativ kann die Stabilität mit der z-Transformation bewertet werden (Kapitel \ref{sec:6.3}).

\begin{example}{Stabiles, grenzstabiles und instabiles System}
\label{beispiel-stabiles-grenzstabiles-und-instabiles-system}

Als Beispiele werden rekursive Systeme untersucht. Das erste System hat die Differenzengleichung
\begin{align}
y_{1}[ k] &= 10  u[ k] + \frac{1}{2}  y_{1}[ k - 1] \label{eq:4.41}
\end{align}
Wird die Anregung \(u[k]\) nach einer endlichen Zeit zu null, wird der Wert \(y_{1}\){[}k{]} nur halb so groß wie der Wert \(y_{1}\){[}k -- 1{]} im Takt zuvor. Aus diesem Grund konvergiert \(y_{1}\){[}k{]} für k → ∞ gegen null. Das System ist asymptotisch stabil. Das zweite System wird über die Differenzengleichung
\begin{align}
y_{2}[ k] &= 2  u[ k] + y_{2}[ k - 1] \label{eq:4.42}
\end{align}
beschrieben. Wird die Anregung \(u[k]\) zu null, bleibt der Wert \(y_{2}\){[}k{]} genauso so groß wie der Wert\\
\(y[k -- 1]\) im Takt zuvor. Aus diesem Grund ist \(y_{2}\){[}k{]} für k → ∞ konstant. Das System ist grenzstabil. Das dritte System wird durch die Gleichung
\begin{align}
y_{3}[ k] &= \frac{1}{2}  u[ k] + \frac{11}{10}  y_{3}[ k - 1] \label{eq:4.43}
\end{align}
definiert. Bereits an der Differenzengleichung wird deutlich, dass auch nach Ende der Anregung der Wert \(y[k]\) einen Faktor 1.1 größer ist als der Werte \(y[k - 1]\) im Takt zuvor. Die Systemantwort divergiert. Das System ist damit instabil. Bild \ref{fig:4.11} zeigt die Sprungantworten der unterschiedlichen Systeme. Alle Systeme werden mit einer Rechteckfolge
\begin{align}
u[ k] &= \sigma[ k] - \sigma[ k - 10] \label{eq:4.44}
\end{align}
angeregt, die eine endliche Energie aufweist.

\begin{fixedfigure}
  \includegraphics[alt={Sprungantwort eines asymptotisch stabilen, grenzstabilen und instabilen Systems}]{new media/4_11_sprungantwort_eines_asymptotisc.png}
  \captionof{figure}{Sprungantwort eines asymptotisch stabilen, grenzstabilen und instabilen Systems}
  \label{fig:4.11}
\end{fixedfigure}

Bei dem asymptotisch stabilen System klingt die Systemantwort ab, wenn das Eingangssignal zu null wird. Bei dem grenzstabilen System bleibt die Systemantwort konstant, wenn die Anregung zu null wird. Bei dem instabilen System wächst das Ausgangssignal stetig an, auch nachdem die Anregung zu null wird.

Die Stabilitätseigenschaften lassen sich bei diesen einfachen Systemen bereits an der Differenzengleichung abgelesen, indem das Eingangssignal zu null gesetzt wird. Bei dem stabilen System ergibt sich der neue Ausgangswert aus einem Bruchteil des alten Ausgangswertes. Bei dem grenzstabilen System sind alter und neuer Ausgangswert identisch. Bei dem instabilen System ergibt sich der neue Ausgangswert aus einem Vielfachen des alten Ausgangswertes. Der Koeffizient vor dem Term \(y[k - 1]\) entscheidet bei diesem System erster Ordnung offenbar über die Stabilität des Systems.

\end{example}

\subsection{\texorpdfstring{Zusammenfassung grundlegender Systemeigenschaften }{Zusammenfassung grundlegender Systemeigenschaften }}\label{zusammenfassung-grundlegender-systemeigenschaften}

Tabelle \ref{tab:4.1} fasst die diskutierten Systemeigenschaften und ihre Bedeutung zusammen.

\begin{summary}[label=tab:4.1]{Zusammenfassung von Systemeigenschaften}
\begin{center}
  \begin{tabular}{ll}
    \toprule
    Eigenschaft & Bedeutung \\
    \midrule
    Linearität & 
System reagiert auf Linearkombination 
von Eingangssignalen

\[u\lbrack k\rbrack = \nu_{1}  u_{1}\lbrack k\rbrack + \nu_{2}  u_{2}\lbrack k\rbrack\]

mit derselben Linearkombination 
von Ausgangssignalen

\[y\lbrack k\rbrack = \nu_{1}  y_{1}\lbrack k\rbrack + \nu_{2}  y_{2}\lbrack k\rbrack\] \\
    Zeitinvarianz & 
System reagiert auf ein verzögertes 
Eingangssignal u{[}k - \(k_{0}\){]} 
mit einem Ausgangsignal y{[}k - \(k_{0}\){]} \\
    Kausalität & 
System reagiert auf ein Eingangssignal 
erst nach Beginn der Anregung 
\(c_{0}\) ≠ 0 \\
    Asymptotische Stabilität & 
System erreicht nach einer Anregung 
mit endlicher Energie wieder seine Ruheposition \\
    Grenzstabilität & 
System bleibt nach einer Anregung 
in der aktuellen Position \\
    Instabilität & System reagiert nach einer Anregung mit endlicher Energie mit einer divergierenden Systemantwort \\
    \bottomrule
  \end{tabular}
\end{center}
\end{summary}

\protect\phantomsection\label{_Ref369028047}{}

\section{Lösung linearer Differenzengleichungen mit konstanten Koeffizienten}\label{sec:4.3}

Die dargestellten Beispiele haben ein Systemverhalten, das über lineare Differenzengleichungen N-ter Ordnung mit konstanten Koeffizienten beschrieben wird. Zur Berechnung des Ausgangssignals im Zeitbereich stehen unterschiedliche Methoden zur Verfügung, die in diesem Kapitel vorgestellt werden.

\subsection{Rekursive Darstellung von Differenzengleichungen}\label{sec:rekursive-darstellung-von-differenzengleichungen}

Die Differenzengleichung N-ter Ordnung mit konstanten Koeffizienten lautet in ihrer allgemeinen Form
\begin{align}
\sum\limits_{n = 0}^{N}{c_{n}  y[ k - n]} &= \sum\limits_{l = 0}^{L}{d_{l}  u[ k - l]} \label{eq:4.45}
\end{align}
Ohne Beschränkung der Allgemeinheit wird die Annahme \(c_{0} = 1\) getroffen. Mit dieser Annahme kann die Gleichung umgeformt werden zu
\begin{align}
y[ k] &= \sum\limits_{l = 0}^{L}{d_{l}  u[ k - l]} - \sum\limits_{n = 1}^{N}{c_{n}  y[ k - n]} \label{eq:4.46}
\end{align}
Das aktuelle Ausgangssignal ergibt sich allgemein aus dem aktuellen Wert des Eingangssignals sowie den vergangenen Werten des Ein- und Ausgangssignals. Die Ausgangssignale der oben dargestellten Systeme werden mit Hilfe dieser rekursiven Darstellung berechnet.

Für numerische Berechnungen reicht dieses Lösungsverfahren für lineare Differenzengleichungen mit konstanten Koeffizienten vollständig aus. Die rekursive Form hat aber den Nachteil, dass keine geschlossene Form des Ausgangssignals angegeben werden kann. Sie ist deshalb nicht geeignet für eine Analyse von Eigenschaften wie Schwingungsneigung und Stabilität. Deshalb werden in diesem Abschnitt Verfahren zur analytischen Berechnung der Systemantwort beschrieben.

\subsection{Explizite Lösung über die Vier-Schritt-Methode}\label{sec:4.3.2}

Für die Diskussion von Systemeigenschaften ist es notwendig, eine geschlossene Darstellung des Ausgangssignals zu erhalten. Für die Lösung von linearen Differenzengleichungen mit konstanten Koeffizienten gibt es eine sogenannte Vier-Schritt-Methode. Sie besteht aus den Schritten

\begin{itemize}
\item
  Berechnung der allgemeinen homogenen Lösung
\item
  Berechnung einer partikulären Lösung
\item
  Kombination von homogener und partikulärer Lösung
\item
  Bestimmung der Konstanten über Anfangsbedingungen
\end{itemize}

\subsubsection{Berechnung der allgemeinen homogenen Lösung}\label{sec:berechnung-der-allgemeinen-homogenen-luxf6sung}

Zur Berechnung der allgemeinen homogenen Lösung wird das Eingangssignal auf null gesetzt. Es ergibt sich die homogene Differenzengleichung
\begin{align}
y_{H}[ k] + \sum\limits_{n = 1}^{N}{c_{n}  y_{H}[ k - n]} &= 0 \label{eq:4.47}
\end{align}
Mit dem Ansatz
\begin{align}
y_{H}[ k] &= Y_{0}  \lambda^{k} \label{eq:4.48}
\end{align}
kann die homogene Lösung gefunden werden. Einsetzen des Ansatzes führt zu der Gleichung
\begin{align}
Y_{0}  \lambda^{k} + Y_{0}  \sum\limits_{n = 1}^{N}{c_{n}  Y_{0}  \lambda^{k - n}} &= 0 \label{eq:4.49}
\end{align}
beziehungsweise
\begin{align}
\lambda^{k} + \sum\limits_{n = 1}^{N}{c_{n}  \lambda^{k - n}} &= 0 \label{eq:4.50}
\end{align}
Sie wird als charakteristische Gleichung bezeichnet, weil mit ihr die für das System charakteristischen Parameter \(λ_{n}\) bestimmt werden. Da es sich bei dem System um ein lineares System handelt, ergibt sich die allgemeine homogene Lösung aus der Linearkombination der berechneten Nullstellen \(λ_{n}\).
\begin{align}
y_{H}[ k] &= Y_{1}  \lambda_{1}^{k} + Y_{2}  \lambda_{2}^{k} + ... + Y_{N}  \lambda_{N}^{k} \label{eq:4.51}
\end{align}
Die Lösungen der charakteristischen Gleichung müssen jedoch nicht die Vielfachheit von eins haben. Existiert ein \(N_{1}\)-facher Wert \(λ_{1}\), ergibt sich die allgemeine homogene Lösung
\begin{align}
y_{H}[ k] &= Y_{1}  \lambda_{1}^{k} + Y_{2}  k  \lambda_{1}^{k} + ... + Y_{N_{1}}  k^{N_{1} - 1}  \lambda_{1}^{k} + Y_{N_{1} + 1}  \lambda_{2}^{k} + ... + Y_{N}  \lambda_{N - N_{1} + 1}^{k} \label{eq:4.52}
\end{align}
\begin{example}{Sprungantwort eines Tiefpasses - Berechnung der allgemeinen homogenen Lösung}
\label{beispiel-sprungantwort-eines-tiefpasses---berechnung-der-allgemeinen-homogenen-luxf6sung}

Das praktische Vorgehen wird am Beispiel des rekursiven Tiefpasses dargestellt. Es soll die Systemreaktion auf einen Sprung am Eingang des Filters berechnet werden. Die Differenzengleichung für den Filter lautet nach Gleichung \eqref{eq:4.1}
\begin{align}
y[ k] &= (1 - GF)  u[ k] + GF  y[ k - 1] \label{eq:4.53}
\end{align}
Für die allgemeine homogene Lösung muss das Eingangssignal \(u[k]\) zu null gesetzt werden.
\begin{align}
y_{H}[ k] - GF  y_{H}[ k - 1] &= 0 \label{eq:4.54}
\end{align}
Mit dem Ansatz
\begin{align}
y_{H}[ k] &= Y_{0}  \lambda^{k} \label{eq:4.55}
\end{align}
ergibt sich
\begin{align}
\lambda^{k} - GF  \lambda^{k - 1} &= 0 \label{eq:4.56}
\end{align}
Die Gleichung ist zum einen für \(λ = 0\) erfüllt. Diese triviale Lösung beschreibt aber das Ausgangssignal, das zu allen Zeiten null ist. Es ist deshalb nicht von Interesse. Die von null verschiedene Nullstelle dieser Gleichung ergibt sich zu
\begin{align}
\lambda &= GF \label{eq:4.57}
\end{align}
Die allgemeine homogene Lösung ergibt sich als Linearkombination der berechneten Lösungen und lautet
\begin{align}
y_{H}[ k] &= Y_{0}  GF^{k} \label{eq:4.58}
\end{align}
Die Konstante \(Y_{0}\) wird in Schritt 4 durch Anfangsbedingungen festgelegt.

\end{example}

\subsubsection{Berechnung einer partikulären Lösung}\label{sec:berechnung-einer-partikuluxe4ren-luxf6sung}

Im nächsten Schritt muss eine partikuläre Lösung der Differenzengleichung
\begin{align}
y_{P}[ k] + \sum\limits_{n = 1}^{N}{c_{n}  y_{P}[ k - n]} &= \sum\limits_{l = 0}^{L}{d_{l}  u[ k - l]} \label{eq:4.59}
\end{align}
bestimmt werden. Als Ansatz für \(y_{P}\){[}k{]} wird ein Signal verwendet, das in Abhängigkeit vom Eingangssignal gewählt wird.

\begin{table}
  \centering
  \caption{Ansatz für partikuläre Lösungen linearer Differenzengleichungen}
  \begin{tabular}{ll}
    \toprule
    Eingangssignal \(u[k]\) & Ansatz für die partikuläre Lösung \(y_{P}\){[}k{]} \\
    \midrule
    Konstante & Konstante \\
    Polynom & Polynom gleichen Grades \\
    Exponentialfolge \(u_{0}^{k}\) & Exponentialfolge ku\textsubscript{0}\textsuperscript{k} \\
    Harmonische Schwingung cos(Ωk) & Acos(Ωk + ϕ) \\
    \bottomrule
  \end{tabular}
\end{table}

Durch Einsetzen dieses Ansatzes in die Differenzengleichung ergibt sich eine partikuläre Lösung \(y_{P}\){[}k{]}. Ist das Eingangssignal \(u[k]\) eine Kombination der dargestellten Eingangssignale, muss als Ansatz für die partikuläre Lösung eine entsprechende Kombination von Ansätzen gewählt werden. Dabei ist es ausreichend, eine beliebige partikuläre Lösung zu finden.

Alternativ kann die partikuläre Lösung durch eine Variation der Konstanten durchgeführt werden. Dabei werden die Konstanten \(Y_{0}\) \ldots{} \(Y_{N}\) der allgemeinen homogenen Lösung als Funktion des Index k variiert. Auflösen nach den variierten Konstanten \(Y_{n}(k)\) führt zu der gesuchten partikulären Lösung.

\begin{example}{Sprungantwort eines Tiefpasses - Berechnung einer partikulären Lösung}
\label{beispiel-sprungantwort-eines-tiefpasses---berechnung-einer-partikuluxe4ren-luxf6sung}

Für das Beispiel des rekursiven Tiefpass-Filters wird die partikuläre Lösung über Variation der Konstanten bestimmt. Der Ansatz für die partikuläre Lösung lautet
\begin{align}
y_{P}[ k] &= Y_{0}(k)  GF^{k} \label{eq:4.60}
\end{align}
Einsetzen in die inhomogene Differentialgleichung führt unter Berücksichtigung des Eingangssignals
\begin{align}
u[ k] &= \sigma[ k] \label{eq:4.61}
\end{align}
zu
\begin{align}
Y_{0}(k)  GF^{k} - GF  Y_{0}(k - 1)  GF^{k - 1} &= (1 - GF)  \sigma[ k] \label{eq:4.62}
\end{align}
Damit kann der Koeffizient \(Y_{0}(k)\) rekursiv beschrieben werden als
\begin{align}
Y_{0}(k) &= Y_{0}(k - 1) + \frac{(1 - GF)}{GF^{k}}  \sigma[ k] \label{eq:4.63}
\end{align}
Da nur eine beliebige partikuläre Lösung erforderlich ist, kann \(Y_{0}\)(- 1) = 0 gesetzt werden. Damit sind auch alle Konstanten \(Y_{0}(k)\) mit k \textless{} 1 null. Für k \textgreater{} - 1 ergibt sich
\begin{align}
Y_{0}(0) &= (1 - GF) \label{eq:4.64} \\
Y_{0}(1) &= (1 - GF) + \frac{(1 - GF)}{GF} \label{eq:4.65}
\end{align}
Eine wiederholte Berechnung führt auf \(Y_{0}(k)\)
\begin{align}
Y_{0}(k) &= \left( 1 + \frac{1}{GF} + \left( \frac{1}{GF} \right)^{2} + ... + \left( \frac{1}{GF} \right)^{k} \right)  (1 - GF)  \sigma[ k] \label{eq:4.66}
\end{align}
Mit der Summenformel für die endliche geometrische Reihe ergibt sich wegen der Definition des Gedächtnisfaktors mit \textbar GF\textbar{} \textless{} 1
\begin{align}
Y_{0}(k) &= \left( \frac{1 - \left( \frac{1}{GF} \right)^{k + 1}}{1 - \frac{1}{GF}} \right)  (1 - GF)  \sigma[ k] \label{eq:4.67}
\end{align}
und die partikuläre Lösung lautet
\begin{align}
y_{P}[ k] &= \left( \frac{1 - \left( \frac{1}{GF} \right)^{k + 1}}{1 - \frac{1}{GF}} \right)  (1 - GF)  \sigma[ k]  GF^{k} = \left( 1 - GF^{k + 1} \right)  \sigma[ k] \label{eq:4.68}
\end{align}
\end{example}

\subsubsection{Kombination von allgemeiner homogener und partikulärer Lösung}\label{sec:kombination-von-allgemeiner-homogener-und-partikuluxe4rer-luxf6sung}

Die komplette Lösung der Differenzengleichung ergibt sich als Summe aus der allgemeinen homogenen und einer partikulären Lösung.
\begin{align}
y[ k] &= Y_{0}  \lambda_{0}^{k} + Y_{1}  \lambda_{1}^{k} + ... + Y_{N}  \lambda_{N}^{k} + y_{P}[ k] \label{eq:4.69}
\end{align}
\begin{example}{Sprungantwort eines Tiefpasses - Kombination der beiden Lösungen}
\label{beispiel-sprungantwort-eines-tiefpasses---kombination-der-beiden-luxf6sungen}

Die Summe von allgemeiner homogener und partikulärer Lösung errechnet sich zu
\begin{align}
y[ k] &= y_{H}[ k] + y_{P}[ k] = Y_{0}  GF^{k} + \left( 1 - GF^{k + 1} \right)  \sigma[ k] \label{eq:4.70}
\end{align}
\end{example}

\subsubsection{Bestimmung der Konstanten über Anfangsbedingungen}\label{sec:bestimmung-der-konstanten-uxfcber-anfangsbedingungen}

Die Werte für \(Y_{n}\) müssen durch Einsetzen von Folgenwerten ermittelt werden. Für jede Konstante \(Y_{n}\) ist ein Folgenwert als Randbedingung notwendig.

\begin{example}{Sprungantwort eines Tiefpasses - Bestimmung der Konstanten}
\label{beispiel-sprungantwort-eines-tiefpasses---bestimmung-der-konstanten}

Für das vorliegende Beispiel wird vorgegeben, dass das Ausgangssignal für k \textless{} 0 null ist. Für\\
\(k = 0\) ergibt sich das Ausgangssignal mit \(y[- 1]\) = 0 direkt über die Differenzengleichung zu
\begin{align}
y[ 0] &= (1 - GF) + GF  0 = 1 - GF \label{eq:4.71}
\end{align}
Der Wert \(Y_{0}\) in Gleichung \eqref{eq:4.70} muss über die Anfangsbedingungen \(y[0]\) = 1 - GF bestimmt werden.
\begin{align}
y[ 0] &= Y_{0}  GF^{0} + 1 - GF^{1} = Y_{0} + 1 - GF = 1 - GF \label{eq:4.72}
\end{align}
Es ergibt sich die Konstante \(Y_{0} = 0\), und die Lösung der Differenzengleichung lautet
\begin{align}
y[ k] &= \left( 1 - GF^{k + 1} \right)  \sigma[ k] \label{eq:4.73}
\end{align}
Bild \ref{fig:4.12} vergleicht das Ausgangssignal der rekursiven Berechnung mit dem der analytischen Berechnung.

\begin{fixedfigure}
  \includegraphics[alt={Vergleich zwischen rekursiver und analytischer Berechnung des Ausgangssignals\\
eines rekursiven Tiefpasses für GF = 0.9}]{new media/4_12_vergleich_zwischen_rekursiver_u.png}
  \captionof{figure}{Vergleich zwischen rekursiver und analytischer Berechnung des Ausgangssignals\\
eines rekursiven Tiefpasses für GF = 0.9}
  \label{fig:4.12}
\end{fixedfigure}

Erwartungsgemäß stimmen die beiden Berechnungen überein.

\end{example}

Das Vorgehen bei der Vier-Schritt-Methode zur Lösung linearer Differenzengleichungen mit konstanten Koeffizienten ist in Tabelle \ref{tab:4.3} zusammengefasst.

\begin{table}
  \centering
  \caption{Vorgehen bei der Vier-Schritt-Methode zur Lösung linearer Differenzengleichungen mit konstanten Koeffizienten}
  \label{tab:4.3}
  \begin{tabular}{ll}
    \toprule
    Schritt & Beschreibung \\
    \midrule
    1 & Lösung der homogenen Differenzengleichung

\(\sum\limits_{n = 0}^{N}{c_{n}  y\lbrack k - n\rbrack} = 0\)

über Ansatz

\(y_{H}\lbrack k\rbrack = Y_{0}  \lambda^{k}\) \\
    & durch Lösen der charakteristischen Gleichung \\
    & \(\lambda^{k} + \sum\limits_{n = 1}^{N}{a_{n}  \lambda^{k - n}} = 0\) \\
    & Allgemeine Lösung in Abhängigkeit der Vielfachheit \\
    & \(y_{H}\lbrack k\rbrack = Y_{1}  \lambda_{1}^{k} + Y_{2}  k  \lambda_{1}^{k} + ... + Y_{P}  k^{P - 1}  \lambda_{1}^{k} + Y_{P + 1}  \lambda_{2}^{k} + ... + Y_{N}  \lambda_{N - P + 1}^{k}\) \\
    2 & Bestimmung einer partikulären Lösung \(y_{P}\){[}k{]}

- über einen Lösungsansatz je nach Eingangssignal

- über die Methode Variation der Konstanten \\
    3 & Superposition von allgemeiner homogener und spezieller partikulärer Lösung \\
    & \(y\lbrack k\rbrack = y_{H}\lbrack k\rbrack + y_{P}\lbrack k\rbrack\) \\
    4 & Bestimmung der nicht definierten Konstanten über Anfangsbedingungen \\
    \bottomrule
  \end{tabular}
\end{table}

\subsection{Stabilität und charakteristische Gleichung eines Systems}\label{sec:4.3.3}

Bei der Einführung des Begriffes der Stabilität in Abschnitt \ref{sec:4.2.3} wird ausgeführt, dass stabile Systeme nach einer Anregung mit endlicher Energie wieder in ihren Ausgangszustand zurückkehren. Das Verhalten des Systems nach der Anregung wird durch die homogene Lösung der Differentialgleichung beschrieben, die in Abschnitt \ref{sec:4.3.2} berechnet wird. Sie setzt sich bei einfachen Lösungen \(λ_{n}\) aus einer Linearkombination von Potenzfolgen zusammen.
\begin{align}
y_{H}[ k] &= Y_{1}  \lambda_{1}^{k} + Y_{2}  \lambda_{2}^{k} + ... + Y_{N}  \lambda_{N}^{k} \label{eq:4.74}
\end{align}
Damit die homogene Lösung zu null wird, müssen alle Lösungen \(λ_{n}\) einen Betrag \textbar \(λ_{n}\)\textbar{} \textless{} 1 aufweisen. Besitzt ein Wert \(λ_{n}\) einen Betrag \textbar \(λ_{n}\)\textbar{} \textgreater{} 1, divergiert der entsprechende Summand aus Gleichung~\eqref{eq:4.74}, und folglich divergiert auch die Lösung der homogenen Differentialgleichung.

Liegt mit \(λ_{1}\) eine P-fache Lösung der charakteristischen Gleichung vor, weisen die zugehörigen Summanden der homogenen Lösung Terme der Form
\begin{align}
y_{H}[ k] &= Y_{1}  \lambda_{1}^{k} + Y_{2}  k  \lambda_{1}^{k} + ... + Y_{P}  k^{P - 1}  \lambda_{1}^{k} + Y_{P + 1}  \lambda_{2}^{k} + ... + Y_{N}  \lambda_{N - P + 1}^{k} \label{eq:4.75}
\end{align}
auf. Da die Exponentialfunktion schneller fällt und wächst als jede Potenz von k, konvergiert diese Summe ebenfalls für einen Betrag \textbar \(λ_{n}\)\textbar{} \textless{} 1, und sie divergiert für einen Betrag \textbar \(λ_{n}\)\textbar{} \textgreater{} 1. Dabei ist es unerheblich, ob die Lösungen \(λ_{n}\) reell oder komplex sind.

Einen Sonderfall stellen Lösungen mit einem Betrag \textbar \(λ_{n}\)\textbar{} = 1 dar.
\begin{align}
y_{H}[ k] &= Y_{1}  1^{k} + Y_{2}  e^{j  \varphi} + Y_{3}  e^{- j  \varphi} + ... \label{eq:4.76}
\end{align}
Die Lösungen sind konstant beziehungsweise schwingen mit konstanter Amplitude. Für den Fall einfacher Lösungen liegt damit weder eine konvergente noch eine divergente Lösung vor. Der Fall entspricht dem diskutierten Fall der Grenzstabilität des zugehörigen Systems.

Besitzt eine Lösung mit einem Betrag \textbar \(λ_{n}\)\textbar{} = 1 eine Vielfachheit von P \textgreater{} 1, entstehen Terme der Form
\begin{align}
{y_{H} \left[k\right]} &=  {Y_{1}  1^{k} +Y_{2}  k 1^{k} +Y_{3}  k^{2}  1^{k} ...} \\
{} & {} & {+Y_{3}  e{}^{j \varphi  k} +Y{}_{4}  k e{}^{j \varphi  k} +Y{}_{5}  k{}^{2}  e{}^{j \varphi  k} ...} \\
{} & {} & {+Y{}_{6}  e{}^{-j \varphi  k} +Y{}_{7}  k e{}^{-j \varphi  k} +Y{}_{8}  k{}^{2}  e{}^{-j \varphi  k} ...}
 \label{eq:4.77}
\end{align}
Da die Exponentialfunktion die Terme nicht dämpft, divergiert der Ausdruck und damit die gesamte homogene Lösung. Das System ist instabil. Aus dieser Diskussion ergibt sich der in Tabelle \ref{tab:4.7} beschriebene Zusammenhang zwischen der Stabilität von linearen, zeitinvarianten Systemen und den Lösungen der charakteristischen Gleichung.

\begin{table}
  \centering
  \caption{Zusammenhang zwischen Lösungen der charakteristischen Gleichung und der Stabilität von LTI-Systemen}
  \begin{tabular}{ll}
    \toprule
    Eigenschaft & Lösungen \(λ_{n}\) der charakteristischen Gleichung \\
    \midrule
    Asymptotisch stabiles System & 
Alle Lösungen \(λ_{n}\) besitzen einen 
Betrag \textbar \(λ_{n}\)\textbar{} \textless{} 1 \\
    Grenzstabiles System & 
Alle Lösungen \(λ_{n}\) besitzen einen 
Betrag \textbar \(λ_{n}\)\textbar{} \textless{} 1, 
zusätzlich liegt mindestens eine 
einfache Lösung mit Betrag \textbar \(λ_{n}\)\textbar{} = 1 vor \\
    Instabiles System & 
Es existiert mindestens eine Lösung \(λ_{n}\) 
mit einem Betrag \textbar \(λ_{n}\)\textbar{} \textgreater{} 1 
oder eine mehrfache Lösung mit Betrag \textbar \(λ_{n}\)\textbar{} = 1 \\
    \bottomrule
  \end{tabular}
\end{table}

\subsection{Sprung- und Impulsantwort eines Systems}\label{sec:4.3.4}

Das Ausgangssignal eines zeitdiskreten Systems ist von dem Anfangszustand abhängig. Sind die Anfangsbedingungen null, ist das System energiefrei. Wie im zeitkontinuierlichen Bereich wird die Reaktion eines energiefreien Systems auf eine sprungförmige Erregung \(σ[k]\) als Sprungantwort \(h[k]\) bezeichnet.

\begin{figure}
  \centering
  \includegraphics[alt={Sprungantwort \(h[k]\) und Impulsantwort \(g[k]\) als Ausgangssignal eines energiefreien Systems}]{new media/4_13_sprungantwort_hk_und_impulsantw.pdf}
  \caption{Sprungantwort \(h[k]\) und Impulsantwort \(g[k]\) als Ausgangssignal eines energiefreien Systems}
\end{figure}

Für das Beispiel des rekursiven Filters wird in Abschnitt \ref{sec:4.3.2} die Antwort des energiefreien Systems auf einen Sprung der Höhe 1 mit der Bedingung \(y[k = - 1]\) = 0 berechnet.
\begin{align}
y[ k] &= \left( 1 - GF^{k + 1} \right)  \sigma[ k] \label{eq:4.78}
\end{align}
Analog wird die Impulsantwort \(g[k]\) als Reaktion eines energiefreien Systems auf eine Anregung mit einem Impuls \(δ[k]\) definiert. Im Kapitel \ref{sec:3.2} wird gezeigt, dass die Impulsfolge \(δ[k]\) als Differenz zweier Sprungfolgen berechnet werden kann.
\begin{align}
\delta[ k] &= \sigma[ k] - \sigma[ k - 1] \label{eq:4.79}
\end{align}
Wegen der Linearität und Zeitinvarianz errechnet sich die Systemreaktion eines energiefreien LTI-Systems auf einen Impuls am Eingang aus der Differenz zweier Sprungantworten zu
\begin{align}
g[ k] &= h[ k] - h[ k - 1] \label{eq:4.80}
\end{align}
Die Systemantwort \(g[k]\) eines rekursiven Filters auf einen Impuls \(δ[k]\) ergibt sich demnach zu
\begin{align}
g[ k] &= h[ k] - h[ k - 1] = \left( 1 - GF^{k + 1}  \sigma[ k] \right) - \left( 1 - GF^{k}  \sigma[ k - 1] \right) \label{eq:4.81}
\end{align}
\subsection{Berechnung der Systemantwort durch Superposition}\label{sec:berechnung-der-systemantwort-durch-superposition}

Ist ein System linear und zeitinvariant, kann ein Ausgangssignal dadurch berechnet werden, dass die Eingangssignale zerlegt, ihre jeweiligen Systemantworten berechnet und anschließend addiert werden. Dieses Prinzip wird als Superpositionsprinzip bezeichnet. Als erste Anwendung dieses Prinzips wird in Abschnitt \ref{sec:4.3.4} die Impulsantwort als Differenz zweier Sprungantworten berechnet. Wird zum Beispiel ein rekursives Filter mit der Differenzengleichung
\begin{align}
y[ k] &= (1 - GF)  u[ k] + GF  y[ k - 1] \label{eq:4.82}
\end{align}
mit einer Rechteckfolge der Länge 10 und der Höhe 5 beaufschlagt, kann das Eingangssignal als Summe zweier Sprungfolgen dargestellt werden
\begin{align}
u[ k] &= 5  \left( \sigma[ k] - \sigma[ k - 10] \right) = 5  \sigma[ k] - 5  \sigma[ k - 10] \label{eq:4.83}
\end{align}
Damit ergibt sich das Ausgangsignal \(y[k]\) aus der Summe der beiden Sprungantworten
\begin{align}
y[ k] &= 5  h[ k] - 5  h[ k - 10] = 5  \left( 1 - GF^{k + 1}  \sigma[ k] \right) - 5  \left( 1 - GF^{k - 9}  \sigma[ k - 10] \right) \label{eq:4.84}
\end{align}
Bild \ref{fig:4.14} stellt das Superpositionsprinzip für das Beispiel des rekursiven Filters bei Anregung mit einem rechteckförmigen Signal dar.

\begin{figure}
  \centering
  \includegraphics[alt={Überlagerung der Systemreaktion \(y[k]\) = \(y_{1}\){[}k{]} + \(y_{2}\){[}k{]} bei überlagertem Eingangssignal \(u[k]\) = \(u_{1}\){[}k{]} + \(u_{2}\){[}k{]}}]{new media/4_14_berlagerung_der_systemreaktion.png}
  \caption{Überlagerung der Systemreaktion \(y[k]\) = \(y_{1}\){[}k{]} + \(y_{2}\){[}k{]} bei überlagertem Eingangssignal \(u[k]\) = \(u_{1}\){[}k{]} + \(u_{2}\){[}k{]}}
  \label{fig:4.14}
\end{figure}

Mit der Kenntnis der Sprungantwort eines Systems kann demnach für Eingangssignale, die sich über die Sprungfolgen darstellen lassen, eine Systemantwort über Superposition berechnet werden.

▪

\section{Berechnung der Systemantwort über die Faltungssumme}\label{sec:4.4}

Das Superpositionsprinzip erlaubt eine Zerlegung des Eingangssignals in eine Linearkombination von Eingangssignalen, deren entsprechendes Ausgangsignal bekannt ist. Auf Basis des Superpositionsprinzips kann bei bekannter Impulsantwort \(g[k]\) das Ausgangssignal zu einem beliebigen Eingangssignal \(u[k]\) bestimmt werden. Der Vorgang wird wie bei zeitkontinuierlichen Systemen als Faltung bezeichnet. Im diskreten Zeitbereich geht das Faltungsintegral in eine sogenannte Faltungssumme über.

\subsection{Herleitung der Faltungssumme}\label{sec:herleitung-der-faltungssumme}

Mathematisch kann die Faltungssumme über das Superpositionsprinzip hergeleitet werden. Die Systemreaktion eines energiefreien Systems auf einen Impuls am Eingang ist die Impulsantwort \(g[k]\). Ein beliebiges Eingangssignal \(u[k]\) kann mit der Ausblendeigenschaft der Impulsfolge als gewichtete Summe von Impulsen beschrieben werden.
\begin{align}
u[ k] &= \sum\limits_{\kappa = - \infty}^{\infty}{u[\kappa]  \delta[ k - \kappa]} \label{eq:4.85}
\end{align}
Die Systemantwort \(y[k]\) auf ein solches Eingangssignal ergibt sich aus derselben Linearkombination von Impulsantworten
\begin{align}
y[ k] &= \sum\limits_{\kappa = - \infty}^{\infty}{u[\kappa]  g[ k - \kappa]} = u[ k]*g[ k] \label{eq:4.86}
\end{align}
Diese Operation wird als Faltungssumme oder auch diskrete Faltung bezeichnet.

\begin{example}{Berechnung der gleitenden Mittelung über die Faltungssumme}
\label{beispiel-berechnung-der-gleitenden-mittelung-uxfcber-die-faltungssumme}

Der im vorangegangenen Abschnitt behandelte Algorithmus zur gleitenden Mittelung führte zu der Differenzengleichung
\begin{align}
y[ k] &= \frac{1}{5}  \left( u[ k] + u[ k - 1] + u[ k - 2] + u[ k - 3] + u[ k - 4] \right) \label{eq:4.87}
\end{align}
Durch Einsetzen der Impulsfolge als Eingangssignal ergibt sich die Impulsantwort zu
\begin{align}
g[ k] &= \frac{1}{5}  \left( \delta[ k] + \delta[ k - 1] + \delta[ k - 2] + \delta[ k - 3] + \delta[ k - 4] \right) \label{eq:4.88}
\end{align}
und das Ausgangssignal zu einem beliebigen Eingangssignal kann durch die Faltungssumme berechnet werden
\begin{align}
{y\left[k\right]} &=  {\sum\limits _{\kappa =-\infty }^{\infty }u\left[\kappa \right] g\left[k-\kappa \right] } \\
{} &=  {\sum\limits _{\kappa =-\infty }^{\infty }u\left[\kappa \right] \frac{1}{5}  \left(\delta \left[k-\kappa \right]+\delta \left[k-\kappa -1\right]+\delta \left[k-\kappa -2\right]+\delta \left[k-\kappa -3\right]+\delta \left[k-\kappa -4\right]\right) } \\
{} &=  {\frac{1}{5}  \sum\limits _{\kappa =0}^{4}u\left[k-\kappa \right] }
 \label{eq:4.89}
\end{align}
\end{example}

\subsection{Grafische Interpretation der Faltungssumme}\label{sec:grafische-interpretation-der-faltungssumme}

Die direkte analytische Berechnung der Faltungssumme ist nur in Ausnahmefällen möglich und sinnvoll, da die Zusammenfassung der Summe aufwendig ist. Im Gegensatz zur zeitkontinuierlichen Faltung kann die zeitdiskrete Faltung aber numerisch ausgeführt werden. Die Faltungssumme ist damit eine Realisierungsform zeitdiskreter Systeme. Zum besseren Verständnis wird deshalb die Faltung an einem Beispiel zweier Folgen grafisch veranschaulicht.

\begin{example}{Grafische Interpretation der Faltungssumme}
\label{beispiel-grafische-interpretation-der-faltungssumme}

Die Folge \(u[k]\) wird als Sprungfolge angenommen, die Impulsantwort \(g[k]\) ergibt sich in diesem Beispiel zu
\begin{align}
g[ k] &= 2  \sigma[ k] - \sigma[ k - 2] - \sigma[ k - 4] \label{eq:4.90}
\end{align}
Die Folgen sind in Bild \ref{fig:4.15} grafisch dargestellt.

\begin{fixedfigure}
  \includegraphics[alt={Folgen \(u[k]\) und \(g[k]\) für das Beispiel zur grafischen Interpretation der Faltungssumme}]{new media/4_15_folgen_uk_und_gk_fr_das_beispie.png}
  \captionof{figure}{Folgen \(u[k]\) und \(g[k]\) für das Beispiel zur grafischen Interpretation der Faltungssumme}
  \label{fig:4.15}
\end{fixedfigure}

Die Faltung ist über eine Summenformel definiert. Sie kann umgeformt werden zu
\begin{align}
u[ k]*g[ k] &= \sum\limits_{\kappa = - \infty}^{\infty}{u[\kappa]  g[ k - \kappa]} = \sum\limits_{\kappa = - \infty}^{\infty}{u[\kappa]  g\left[ - (\kappa - k) \right]} \label{eq:4.91}
\end{align}
Mit dieser Darstellung fällt die Interpretation einfacher: Die Folge \(g[κ]\) wird an der Achse \(κ = 0\) gespiegelt und um k nach rechts verschoben. Dann wird das Produkt der einzelnen Folgenwerte addiert. Bild \ref{fig:4.16} stellt die unterschiedlichen Phasen der grafischen Faltung dar.

\begin{fixedfigure}
  \includegraphics[alt={Grafische Veranschaulichung der Faltung zweier Folgen zu verschiedenen Stadien}]{new media/4_16_grafische_veranschaulichung_der.png}
  \captionof{figure}{Grafische Veranschaulichung der Faltung zweier Folgen zu verschiedenen Stadien}
  \label{fig:4.16}
\end{fixedfigure}

Für negative Folgenindizes k überschneiden sich die beiden Folgen nicht. Zum Zeitpunkt \(k = 0\) überschneiden sich die beiden Folgen an genau einer Stelle \(κ = 0\). Das Ergebnis ist damit \(y[0]\) = 2. Für\\
\(k = 1\) ergibt sich eine Überschneidung der ersten beiden Werte. Nach Bildung des Produktes werden die Ergebnisse addiert und es ergibt sich \(y[1]\)= 2 + 2 = 4. So wird für die übrigen Werte von k fortgefahren. Für k ≥ 3 überschneidet sich die Folge x komplett mit der Folge g, sodass sich der Wert des Ausgangssignals nicht weiter ändert.

\begin{fixedfigure}
  \includegraphics[alt={Ergebnis der Faltung zweier Folgen}]{new media/4_17_ergebnis_der_faltung_zweier_fol.png}
  \captionof{figure}{Ergebnis der Faltung zweier Folgen}
\end{fixedfigure}

Aus der Grafik kann abgelesen werden, dass das Signal \(y[k]\) ab \(k = 3\) konstant bleibt.

\end{example}

Die in diesem Beispiel dargestellte Methode zur Berechnung der Faltungssumme ist in Tabelle \ref{tab:4.5} zusammengefasst.

\begin{table}
  \centering
  \caption{Vorgehen bei der Berechnung der Systemantwort über die Faltungssumme}
  \label{tab:4.5}
  \begin{tabular}{ll}
    \toprule
    Schritt & Beschreibung \\
    \midrule
    1 & Berechnung der Impulsantwort \(g[k]\) \\
    2 & Skizze von Eingangssignal \(u[κ]\) und Impulsantwort \(g[κ]\) \\
    3 & 
Skizze von einem der Signale \(u[k - κ]\) oder \(g[k - κ]\) über Spiegelung an der Achse \(κ = 0\) und 
Verschiebung um k nach rechts \\
    4 & 
Aufteilen der Faltungssumme in sinnvolle Bereiche 
(Überlappungsbereiche, Sprungstellen, Definitionsgrenzen, ...) \\
    5 & Lösen der Summen und Superposition der Ergebnisse \\
    \bottomrule
  \end{tabular}
\end{table}
\begin{application}{Zeitdiskrete Faltung}
  Im Online-Portal \emph{Systemtheorie Online} verdeutlicht die Applikation \emph{Zeitdiskrete Faltung} grafisch die Faltungsoperation für zeitdiskrete Signale.
  \tcblower
  \url{h-ka.de/eit/digitale-signalverarbeitung/übungen/kapitel2}
\end{application}
\subsection{Rechenregeln zur Faltungssumme}\label{sec:rechenregeln-zur-faltungssumme}

Zur Vereinfachung der Berechnung von Faltungssummen können Rechenregeln angewendet werden, die im Folgenden kurz dargestellt sind. Aus der grafischen Darstellung zur Faltung wird deutlich, dass die Faltung eine kommutative Operation ist. Der mathematische Nachweis ergibt sich aus einer Indextransformation.
\begin{align}
x_{1}[ k]*x_{2}[ k] &= \sum\limits_{\kappa = - \infty}^{\infty}{x_{1}[\kappa]  x_{2}[ k - \kappa]} = \sum\limits_{k - \kappa = - \infty}^{\infty}{x_{1}[ k - \kappa]  x_{2}[\kappa]} \label{eq:4.92}
\end{align}
Da sich die Summe von - ∞ bis + ∞ erstreckt, ist eine Verschiebung um k und eine Spiegelung nicht relevant, und es ergibt sich
\begin{align}
x_{1}[ k]*x_{2}[ k] &= \sum\limits_{\kappa = - \infty}^{\infty}{x_{1}[\kappa]  x_{2}[ k - \kappa]} = \sum\limits_{\kappa = - \infty}^{\infty}{x_{1}[ k - \kappa]  x_{2}[\kappa]} = x_{2}[ k]*x_{1}[ k] \label{eq:4.93}
\end{align}
Das Distributivgesetz ergibt sich aus der Linearität der Summe.
\begin{align}
{\left( x_{1}[ k] + x_{2}[ k] \right)*x_{3}[ k] = \sum\limits_{\kappa = - \infty}^{\infty}{\left( x_{1}[\kappa] + x_{2}[\kappa] \right)  x_{3}[ k - \kappa]}
}{= \sum\limits_{\kappa = - \infty}^{\infty}{x_{1}[\kappa]  x_{3}[ k - \kappa]} + x_{2}[\kappa]  x_{3}[ k - \kappa] = x_{1}[ k]*x_{3}[ k] + x_{2}[ k]*x_{3}[ k]} \label{eq:4.94}
\end{align}
Das Assoziativgesetz wird hier nur genannt und nicht bewiesen. Es lautet:
\begin{align}
\left( x_{1}[ k]*x_{2}[ k] \right)*x_{3}[ k] &= x_{1}[ k]*\left( x_{2}[ k]*x_{3}[ k] \right) \label{eq:4.95}
\end{align}
Bei der grafischen Faltung wird gezeigt, dass der von null verschiedene Überlappungsbereich für kausale Folge immer im Zahlenbereich von 0 \ldots{} k liegt. Aus diesem Grund muss die Faltung auch nur in diesem Bereich ausgeführt werden, sodass für die Faltung zweier kausaler Folgen gilt:
\begin{align}
\sum\limits_{\kappa = - \infty}^{\infty}{x_{1}[\kappa]  x_{2}[ k - \kappa]} &= \sum\limits_{\kappa = 0}^{k}{x_{1}[\kappa]  x_{2}[ k - \kappa]} \label{eq:4.96}
\end{align}
Aus der Eigenschaft folgt außerdem, dass die Faltung zweier kausaler Folgen zu einer kausalen Folge führt. Die Rechenregeln für die Faltung mit einer Impulsfolge werden in Abschnitt \ref{sec:3.2.1} bei der Einführung von Impulsfolgen behandelt.

\begin{summary}[label=]{Zusammenfassung der Rechenregeln zur Faltungssumme}
\begin{center}
  \begin{tabular}{ll}
    \toprule
    Rechenregel & Darstellung als Gleichung \\
    \midrule
    Kommutativgesetz & \(x_{1}\lbrack k\rbrack*x_{2}\lbrack k\rbrack = x_{2}\lbrack k\rbrack*x_{1}\lbrack k\rbrack\) \\
    Distributivgesetz & \(\left( x_{1}\lbrack k\rbrack + x_{2}\lbrack k\rbrack \right)*x_{3}\lbrack k\rbrack = x_{1}\lbrack k\rbrack*x_{3}\lbrack k\rbrack + x_{2}\lbrack k\rbrack*x_{3}\lbrack k\rbrack\) \\
    Assoziativgesetz & \(\left( x_{1}\lbrack k\rbrack*x_{2}\lbrack k\rbrack \right)*x_{3}\lbrack k\rbrack = x_{1}\lbrack k\rbrack*\left( x_{2}\lbrack k\rbrack*x_{3}\lbrack k\rbrack \right)\) \\
    Faltung kausaler Folgen 
führt zu einer kausalen Folge
 & \(\sum\limits_{\kappa = - \infty}^{\infty}{x_{1}\lbrack\kappa\rbrack  x_{2}\lbrack k - \kappa\rbrack} = \sum\limits_{\kappa = 0}^{k}{x_{1}\lbrack\kappa\rbrack  x_{2}\lbrack k - \kappa\rbrack}\) \\
    Faltung mit einem Impuls & \(\delta\lbrack k\rbrack*x\lbrack k\rbrack = x\lbrack k\rbrack\) \\
    Faltung mit einem Impuls an der Stelle \(k_{0}\) & \(\delta\left\lbrack k - k_{0} \right\rbrack*x\lbrack k\rbrack = x\left\lbrack k - k_{0} \right\rbrack\) \\
    \bottomrule
  \end{tabular}
\end{center}
\end{summary}

\subsection{Impulsantwort und Stabilität}\label{sec:4.4.4}

In Abschnitt \ref{sec:4.2.3} wird die Stabilität von Systemen aus physikalischer Sicht definiert. Mit dem Wissen, dass sich bei einem LTI-System die Systemantwort \(y[k]\) aus dem Faltungsintegral ergibt, kann die Stabilitätsbewertung auf die Impulsantwort \(g[k]\) zurückgeführt werden. Dabei wird davon ausgegangen, dass das System für den Zeitraum 0 \textless{} k ≤ \(k_{0}\) angeregt wird. Für den Zeitraum k ≥ \(k_{0}\) nach der Anregung wird das Verhalten der Systemantwort \(y[k]\) analysiert.
\begin{align}
y[ k] &= \sum\limits_{\kappa = 0}^{k_{0}}{u[\kappa]  g[ k - \kappa]} \label{eq:4.97}
\end{align}
Aus der physikalischen Bedingung an Stabilität leitet sich die Forderung ab, dass bei einer zeitlich begrenzten Anregung das Ausgangssignal den Grenzwert
\begin{align}
\lim_{k \rightarrow \infty}\left| y[ k] \right| &= 0 \label{eq:4.98}
\end{align}
aufweisen muss. Ist der Betrag des Eingangssignals beschränkt, kann er mit \textbar \(u[κ]\)\textbar{} \textless{} \(u_{MAX}\) abgeschätzt werden, und der Betrag des Ausgangssignals kann abgeschätzt werden mit
\begin{align}
\left| y[ k] \right| \leq \sum\limits_{\kappa = 0}^{k_{0}}{\left| u[\kappa] \right|  \left| g[ k - \kappa] \right|} \leq u_{MAX}  \sum\limits_{\kappa = 0}^{k_{0}}\left| g[ k - \kappa] \right| \label{eq:4.99}
\end{align}
Es handelt sich um eine Summe von \(k_{0}\) + 1 Folgegliedern. Die Summe wird zu null, wenn der Betrag der Impulsantwort \(g[k]\) gegen null konvergiert.
\begin{align}
\lim_{k \rightarrow \infty}\left| g[ k] \right| &= 0 \label{eq:4.100}
\end{align}
Ein System ist damit stabil, wenn die Impulsantwort gegen null konvergiert, es ist instabil, wenn die Impulsantwort divergiert. Einen Sonderfall stellen Impulsantworten \(g[k]\) dar, die für k → ∞ einem konstanten Wert \(g_{0}\) zustreben. Bei diesen Systemen konvergiert das Ausgangssignal für k ≥ \(k_{0}\) gegen einen konstanten Wert.
\begin{align}
\lim_{k \rightarrow \infty}y[ k] &= \lim_{k \rightarrow \infty}\sum\limits_{\kappa = 0}^{k_{0}}{u[\kappa]  g[ k - \kappa]} = g_{0}  \sum\limits_{\kappa = 0}^{k_{0}}{x[\kappa]} = y_{0} \label{eq:4.101}
\end{align}
Systeme, deren Impulsantworten \(g[k]\) für k → ∞ einem konstanten Wert \(g_{0}\) zustreben, entsprechen damit den Bedingungen grenzstabiler Systeme. Dasselbe gilt für Systeme, deren Impulsantwort für\\
k → ∞ mit konstanter Amplitude schwingt. Der Zusammenhang zwischen Impulsantwort und Stabilität linearer, zeitinvarianter Systeme ist in Tabelle \ref{tab:4.7} zusammengefasst.

\begin{summary}[label=tab:4.7]{Zusammenfassung des Zusammenhangs zwischen Impulsantwort und Stabilität von LTI-Systemen}
\begin{center}
  \begin{tabular}{ll}
    \toprule
    Eigenschaft & Bedeutung \\
    \midrule
    Stabiles System & \(\lim_{k \rightarrow \infty}g\lbrack k\rbrack = 0\) \\
    Grenzstabiles System & 
\[\lim_{k \rightarrow \infty}g\lbrack k\rbrack = g_{0}\]

oder harmonische Schwingung 
mit konstanter Amplitude \\
    Instabiles System & \(\lim_{k \rightarrow \infty}g\lbrack k\rbrack\text{ ist divergent}\) \\
    \bottomrule
  \end{tabular}
\end{center}
\end{summary}

Zur Stabilitätsbewertung von Systemen im Zeitbereich muss die Impulsantwort bekannt sein. Es wird sich zeigen, dass eine Bewertung der Stabilität im sogenannten z-Bereich praktikabler vorgenommen werden kann.
